\chapter{Curriculum vitae}
\label{ch:cv}

Ce premier chapitre dresse le bilan de mon activité depuis la thèse : parcours,
recherche, encadrement, rayonnement, enseignement et responsabilités
collectives. Il se termine par la liste classée de ma production scientifique et
par deux annexes détaillant le service d'enseignement et les encadrements
doctoraux.



% ============================================================================
\section{Identification et parcours}
\label{sec:cv-identification}
% ============================================================================

\begin{cvfields}
  \item[Nom] Ammar Mian
  \item[Corps] Maître de conférences, classe normale, échelon 5
  \item[Section \textsc{cnu}] 61 --- Génie informatique, automatique et
    traitement du signal
  \item[Établissement] Université Savoie Mont-Blanc, Polytech Annecy-Chambéry
  \item[Laboratoire] \textsc{listic} (\textsc{ea}~3703)
  \item[Contact] \href{mailto:ammar.mian@univ-smb.fr}{ammar.mian@univ-smb.fr}
    \quad ---\quad \url{https://ammarmian.fr}
\end{cvfields}


\subsection{Parcours professionnel}
\label{subsec:cv-carriere}

\cventry[à partir de 2026]{Délégation \textsc{cnrs}}{\textsc{gipsa}-lab,
  équipe \textsc{gaia}, Grenoble}
À compter de septembre 2026. Cette délégation prolonge mes travaux sur la
géométrie des données multivariées vers la polarimétrie et l'interférométrie
radar, en collaboration avec Nicolas Le Bihan et Pierre-Olivier Amblard.

\cventry[depuis 2020]{Maître de conférences}{Université Savoie Mont-Blanc ---
  \textsc{listic} et Polytech Annecy-Chambéry}
Recherche au \textsc{listic}, enseignement à Polytech Annecy-Chambéry.
Responsable du thème AFUTE (Apprentissage, Fusion, Télédétection)
depuis mars 2024 ; chargé de mission « intelligence artificielle dans les
formations » depuis 2023.

\cventry[2019--2020]{Post-doctorant}{Université Aalto, Helsinki, Finlande}
\emph{Géométrie riemannienne pour l'apprentissage statistique}, sous la
responsabilité du \textsc{pr}~Esa Ollila.

\cventry[2016--2019]{Doctorant}{Laboratoire \textsc{sondra}
  (CentraleSupélec) et université Paris-Saclay}
\emph{Contributions à l'analyse de séries temporelles d'images
\textsc{sar}}~\cvcite{THESE19}.

\cventry[2016]{Enseignant vacataire}{\textsc{iut} d'Orsay, département informatique}
Programmation C++ et Java, projet de robotique.


\subsection{Formation}
\label{subsec:cv-formation}

\cventry[2019]{Doctorat de l'université Paris-Saclay}{Spécialité traitement du
  signal, préparé au laboratoire \textsc{sondra} (CentraleSupélec)}

\begin{cvfields}
  \item[Titre] \emph{Contributions to \textsc{sar} Image Time Series
    Analysis}~\cvcite{THESE19}
  \item[Soutenance] 26 septembre 2019
  \item[Direction] Guillaume Ginolhac et Jean-Philippe Ovarlez ;
    encadrement : Abdourahmane Atto
  \item[Jury]     
    \begin{itemize}[label=$\bullet$]
        \item Jean-Yves Tourneret, Professeur INP-ENSEIHHT  (rapporteur)
        \item André Ferrari, PU Université de Nice (rapporteur)
        \item Frédéric Pascal, Professeur CentraleSupélec (président)
        \item Maria Greco, PU Université de Pise (Examinatrice)
    \end{itemize}
  \item[Distinction] Prix de la meilleure contribution scientifique, école
    doctorale \textsc{stic}, 2019
\end{cvfields}

\cventry[2016]{Diplôme d'ingénieur}{Grenoble \textsc{inp} --- Phelma, spécialité
  traitement du signal}

\cvrubric{Formation pédagogique}

J'ai suivi le cycle de formation à la pédagogie pour l'enseignement supérieur du
département \textsc{apprendre} de l'université Savoie Mont-Blanc, soit 32~h
équivalent \textsc{td}. Dans ce cadre, j'ai développé pour le module
\textsc{info}501 une évaluation automatisée des travaux pratiques qui donne aux
étudiants un retour immédiat et les fait progresser par itérations successives ;
ce développement a été reconnu par une décharge de 32~h équivalent \textsc{td}.


% ============================================================================
\section{Activités de recherche}
\label{sec:cv-recherche}
% ============================================================================

\subsection{Thématiques}
\label{subsec:cv-thematiques}

\begin{cvfields}
  \item[Mots-clés] Télédétection, imagerie satellitaire, apprentissage
    statistique, géométrie riemannienne, statistiques robustes, optimisation
    sur variétés, science reproductible
\end{cvfields}

\medskip

Ma recherche articule deux versants complémentaires : le développement théorique
d'estimateurs et de détecteurs d'une part, la mise en pratique d'algorithmes
rapides et robustes pour de grands volumes de données de télédétection d'autre
part. Le fil conducteur en est la matrice de covariance --- objet de ce mémoire
--- prise comme descripteur des données, et la géométrie riemannienne de
l'espace des matrices symétriques définies positives comme cadre naturel pour
l'estimer et l'exploiter. Mes contributions se répartissent sur quatre axes.

\begin{cvfields}
  \item[Télédétection] Traitement et classification d'images satellitaires,
    radar et géophysiques : séries temporelles \textsc{sar}, détection de
    changements, classification d'objets enfouis à partir de radargrammes
    \textsc{gpr}, analyse de texture par ondelettes adaptées. C'est le versant
    applicatif de mes travaux, mené avec des partenaires
    industriels~\cvcite{JSTARS25,TGRS23,TGRS19,JSTARS20,CHAP21a,CHAP22,ICASSP24,IGARSS24}.
  \item[Géométrie riemannienne] Estimation et apprentissage sur la variété des
    matrices symétriques définies positives : moyenne de Fréchet corrigée en
    grande dimension, lois sur variétés et estimation robuste, apprentissage de
    modèles graphiques factoriels, structures de Kronecker, et plus récemment
    apprentissage profond sur variété
    \textsc{sdp}~\cvcite{ICML24,SP20,EUSIPCO21,ECML23,GRETSI23c,EUSIPCO22,IGARSS22,EUSIPCO26a}.
  \item[Estimation et détection] Détection de cibles multivariées pour réseaux
    d'antennes sous hypothèse non gaussienne, en imagerie \textsc{sar} puis en
    milieu sous-marin --- sonar à croix de
    Mills~\cvcite{TSP19,GRSL19,JOE24,SP24,JOE26}.
  \item[Frugalité] Outils de gestion et d'archivage d'expériences
    scientifiques, mesure de la frugalité des algorithmes, apprentissage
    fédéré décentralisé~\cvcite{EUSIPCO25,MEAS25,IGARSS25}.
\end{cvfields}


\subsection{Cinq publications significatives}
\label{subsec:cv-publications-significatives}

\cventry[JSTARS25]{Classification of Buried Objects From Ground Penetrating Radar Images
  by Using Second-Order Deep Learning Models}{\underline{D.~Jafuno},
  \textbf{A.~Mian}, G.~Ginolhac, N.~Stelzenmuller. \emph{\textsc{ieee} Journal
  of Selected Topics in Applied Earth Observations and Remote Sensing}, 2025}
Aboutissement de la thèse de Douba Jafuno. L'article montre qu'un modèle profond
du second ordre --- où la couche de représentation produit une matrice de
covariance traitée dans sa géométrie propre --- classe les objets enfouis à
partir de radargrammes \textsc{gpr} plus finement qu'une architecture
convolutive standard, sur des données industrielles acquises en conditions
réelles.

\cventry[ICML24]{Random Matrix Theory Improved Fréchet Mean of Symmetric Positive
  Definite Matrices}{F.~Bouchard, \textbf{A.~Mian}, M.~Tiomoko, G.~Ginolhac,
  F.~Pascal. \emph{\textsc{icml}}, Vienne, 2024. Conférence de rang
  \textsc{a}$^\ast$}
Contribution théorique : la moyenne de Fréchet, estimateur central de tout
traitement riemannien de matrices \textsc{sdp}, est biaisée en grande dimension
lorsque le nombre d'échantillons est du même ordre que la dimension. Nous en
donnons une correction issue de la théorie des matrices aléatoires, et montrons
le gain qu'elle procure sur des tâches de classification.

\cventry[JOE24]{Robust Detection for Mills Cross Sonar}{\underline{O.~Lerda},
  \textbf{A.~Mian}, G.~Ginolhac, J.-P.~Ovarlez, D.~Charlot.
  \emph{\textsc{ieee} Journal of Oceanic Engineering}, 2024}
Issu de la thèse d'Olivier Lerda, en partenariat avec iXBlue. Transposition au
sonar à croix de Mills des détecteurs robustes développés en radar : la prise en
compte du caractère non gaussien du fouillis sous-marin et de la structure de la
matrice de covariance améliore nettement la détection lorsque les données
secondaires sont peu nombreuses.

\cventry[ECML23]{Learning Graphical Factor Models with Riemannian
  Optimization}{A.~Hippert-Ferrer, F.~Bouchard, \textbf{A.~Mian}, T.~Vayer,
  A.~Breloy. \emph{\textsc{ecml-pkdd}}, Turin, 2023. Conférence de rang
  \textsc{a}}
Apprentissage de modèles graphiques dont la matrice de précision se décompose en
un terme de faible rang et un terme structuré. La formulation riemannienne
permet d'optimiser conjointement les deux composantes sous contrainte de définie
positivité.

\cventry[TGRS23]{New Robust Sparse Convolutional Coding Inversion Algorithm for Ground
  Penetrating Radar Images}{\underline{M.~Gallet}, \textbf{A.~Mian},
  G.~Ginolhac, E.~Ollila, N.~Stelzenmuller. \emph{\textsc{ieee} Transactions on
  Geoscience and Remote Sensing}, 2023}
Issu du stage de master~2 de Matthieu Gallet. Nouvel algorithme d'inversion pour le
codage convolutif parcimonieux appliqué au \textsc{gpr}, robuste aux aberrations
et aux hétérogénéités du sous-sol.


\subsection{Encadrement doctoral et scientifique}
\label{subsec:cv-encadrement}

À ce jour j'ai co-encadré deux thèses soutenues et j'encadre deux thèses en
cours, auxquelles s'ajoutera celle de Roman Wild à partir de 2026, soit un taux
d'encadrement cumulé de \textbf{204~\%}. S'y ajoutent deux post-doctorants, un
ingénieur de recherche et sept stagiaires de master~2. Le détail de chaque
encadrement --- financement, jury, devenir professionnel --- figure en
\cref{sec:cv-annexe-encadrements}.

\cvrubric{Thèses soutenues}

\cventry[2022--2025]{Douba Jafuno}{Classification de signaux \textsc{gpr} ---
  \textsc{cifre} Géolithe}
\begin{cvfields}
  \item[Soutenance] 4 novembre 2025
  \item[Encadrement] 45~\%, avec Guillaume Ginolhac et Nickolas Stelzenmuller
  \item[Publications] \cvcite{JSTARS25,IGARSS24,GRETSI23b}
\end{cvfields}

\cventry[2021--2024]{Olivier Lerda}{Détection robuste pour un sonar à croix de
  Mills --- ingénieur iXBlue à mi-temps}
\begin{cvfields}
  \item[Calendrier] Début en avril 2021, soutenance le 17 décembre 2024
  \item[Encadrement] 33~\%, avec Guillaume Ginolhac et Jean-Philippe Ovarlez
  \item[Particularité] Thèse menée à mi-temps, en parallèle de son poste
    d'ingénieur chez iXBlue
  \item[Publications] \cvcite{JOE24,JOE26,GRETSI23a}
  \item[Devenir] Ingénieur de recherche chez Exail
\end{cvfields}

\cvrubric{Thèses en cours}

\cventry[depuis 2024]{Matthieu Verlynde}{Apprentissage frugal pour la télédétection ---
  bourse ministérielle}
\begin{cvfields}
  \item[Calendrier] Début octobre 2024, soutenance prévue en 2027
  \item[Encadrement] 60~\%, avec Yajing Yan
  \item[Publications] \cvcite{IGARSS25,EUSIPCO25}
\end{cvfields}

\cventry[depuis 2025]{Wenli Shi}{Réseaux de diffusion parcimonieux sur données non
  textuelles --- \textsc{cifre} Givaudan}
\begin{cvfields}
  \item[Calendrier] Début décembre 2025, soutenance prévue en 2028
  \item[Encadrement] Co-direction à 33~\%, avec Frédéric Pascal (\textsc{l2s},
    université Paris-Saclay) et Matthieu Jonckheere (\textsc{dr cnrs},
    \textsc{laas})
\end{cvfields}

\cventry[à partir de 2026]{Roman Wild}{Phase géométrique en \textsc{polinsar}
  --- chaire \textsc{miai}}
\begin{cvfields}
  \item[Calendrier] Début en octobre 2026
  \item[Encadrement] 33~\%, avec Nicolas Le Bihan et Pierre-Olivier Amblard
  \item[Laboratoire] \textsc{gipsa}-lab, équipe \textsc{gaia}
\end{cvfields}

\cvrubric{Post-doctorants}

J'encadre deux post-doctorants accueillis au laboratoire \textsc{l2s}
(CentraleSupélec).

\cventry{Hugo Brehier}{Mécanismes d'attention et transformers pour données sur
  la variété \textsc{sdp}}
Avec Florent Bouchard et Sylvain Chevallier entre Octobre 2024 et Décembre 2025.

\cventry{Thibault Pautrel}{Apprentissage fédéré par géométrie riemannienne}
Avec Florent Bouchard et Guillaume Ginolhac avec un début Septembre 2025 pour une fin Aout 2027.

\cvrubric{Ingénieur de recherche}

Dans le cadre de ma mission « intelligence artificielle », j'encadre un
ingénieur de recherche dédié au développement d'une plateforme d'\textsc{ia}
souveraine hébergée par l'école.

\cvrubric{Stages de master 2}

\cventry[2026]{Robin Duss}{Phase géométrique en \textsc{polinsar}}
Avril à septembre 2026, encadrement à 33~\%, avec Nicolas Le Bihan et
Pierre-Olivier Amblard au \textsc{gipsa}-lab. Stage préparatoire à la thèse de
R.~Wild sur le même sujet.

\cventry[2025]{Thibault Pautrel}{Apprentissage fédéré par géométrie riemannienne}
Encadrement à 33~\%, avec F.~Bouchard et G.~Ginolhac. Devenu post-doctorant.

\cventry[2025]{Mamour Ndiaye}{Pansharpening par variation totale}
Encadrement à 33~\%, avec Y.~Mhiri et A.~Bhanot.

\cventry[2025]{Thomas Petellat}{Détection de fissures volcaniques par apprentissage
  automatique}
Encadrement à 33~\%, avec C.~Lin-Kwong-Chon et Y.~Yan. Publication :
\cvcite{RSV25}.

\cventry[2024]{Matthieu Verlynde}{Apprentissage frugal en télédétection}
Encadrement à 100~\%. Devenu doctorant sous ma direction.

\cventry[2024]{Emma Molière}{Apprentissage frugal en télédétection}
Encadrement à 50~\% avec Argheesh Bhanot. Devenue doctorante.

\cventry[2020]{Matthieu Gallet}{Inversion robuste de signaux \textsc{gpr}}
Encadrement à 50~\%, avec G.~Ginolhac. Devenu doctorant ; publication
associée : \cvcite{TGRS23}.

\label{subsec:cv-projets}

\cvrubric{Appels à projets financés}

\cventry[2025]{ATROCE --- Apprentissage fédéré décentralisé et frugal sur
  des données continues}{\textsc{aap} Recherche \textsc{usmb},
  10~k\texteuro{} --- \underline{co-porteur}}
Avec S.~Monnet, A.~Benoit, M.~Bettineli, S.~Plassart, G.~Ginolhac.

\cventry[2024]{AFREU --- Apprentissage fédéré riemannien résilient et
  dynamique}{\textsc{aap} Recherche \textsc{usmb}, 15~k\texteuro{} ---
  \underline{co-porteur}}
Avec S.~Monnet, A.~Benoit, M.~Bettineli, S.~Plassart, G.~Ginolhac.

\cventry[2021]{ELABORATE --- Plateforme ouverte pour la classification
  géométrique robuste de signaux \textsc{eeg}}{\textsc{aap} Recherche
  \textsc{usmb}, 3~k\texteuro{}, et appel exploratoire \textsc{gdr-iasis},
  7~k\texteuro{} --- \underline{co-porteur} des deux}
Avec F.~Bouchard et S.~Fortunati. Publication associée : \cvcite{EUSIPCO22}.

\cventry[2020]{SMGA --- Stratigraphie de montagne aéroportée}{\textsc{aap}
  Recherche et montagne \textsc{usmb}, 11~k\texteuro{} ---
  \underline{porteur}}
Avec G.~Ginolhac, A.~Atto et E.~Trouvé. Algorithmes d'inversion de données
\textsc{gpr} géophysiques.

\cvrubric{Contrats de gré à gré}

\cventry[2025--2028]{Accompagnement de thèse \textsc{cifre}, CentraleSupélec}{25~k\texteuro{}
  --- \underline{porteur}}
Réseaux de diffusion parcimonieux, thèse de W.~Shi.

\cventry[2022--2025]{Contrat \textsc{cifre}, Géolithe}{24~k\texteuro{} ---
  \underline{porteur}}
Stratigraphie de montagne aéroportée, thèse de D.~Jafuno.

\cventry[2021--2024]{Contrat \textsc{cifre}, iXBlue}{30~k\texteuro{} ---
  \underline{participant}}
Détection robuste pour sonar à croix de Mills, thèse d'O.~Lerda.

\cventry[2020--2021]{Contrat Géolithe}{12~k\texteuro{} --- \underline{participant}}
Apprentissage automatique pour le traitement de données géoradar aéroportées.


\subsection{Collaborations}
\label{subsec:cv-collaborations}

\begin{cvfields}
  \item[Aalto] Université Aalto, Finlande (Esa Ollila). Collaboration
    entretenue depuis mon post-doctorat par des séjours réciproques : j'ai
    invité E.~Ollila un mois en 2023 et une semaine en 2024, et effectué
    plusieurs séjours de recherche à Helsinki. Publications :
    \cvcite{SP26,TGRS23}.
  \item[\textsc{l2s}] CentraleSupélec (Frédéric Pascal, Florent Bouchard).
    Co-encadrement des deux post-doctorants et de la thèse de W.~Shi, projets
    communs en géométrie riemannienne : \cvcite{ICML24,ECML23,CHAP21a,CHAP22}.
  \item[\textsc{sondra}] Université Paris-Saclay (Jean-Philippe Ovarlez).
    Collaboration continue depuis la thèse, en détection robuste radar et
    sonar : \cvcite{TSP19,JOE24,JOE26}.
  \item[\textsc{cedric}] \textsc{cnam} (Arnaud Breloy). Apprentissage de
    modèles graphiques et optimisation riemannienne, ainsi que l'estimation en
    ligne sous structure de Kronecker :
    \cvcite{ECML23,GRETSI23c,SP24,CHAP22}.
  \item[\textsc{ign}] Alexandre Hippert-Ferrer. Apprentissage de graphes
    structurés et classification riemannienne en présence de données
    manquantes : \cvcite{ECML23,GRETSI23c,EUSIPCO22}.
  \item[Industriels] Géolithe (\textsc{gpr}), iXBlue (sonar), Givaudan
    (réseaux de diffusion).
\end{cvfields}


% ============================================================================
\section{Rayonnement scientifique}
\label{sec:cv-rayonnement}
% ============================================================================

\subsection{Animation et responsabilités scientifiques}
\label{subsec:cv-animation}

\cventry[depuis 2024]{Responsable du thème AFUTE}{Apprentissage, Fusion,
  Télédétection --- \textsc{listic}, 17 permanents et 8 doctorants}
J'organise les réunions et séminaires de l'équipe, définis ses orientations
scientifiques et ai mis en place ses outils de travail collaboratif.

\cventry[2022--2024]{Co-animateur de l'action « Télédétection »}{\textsc{gdr-iasis},
  groupement de recherche \textsc{cnrs}}
Le \textsc{gdr-iasis} fédère la communauté nationale du traitement du signal et
des images ; sa restructuration en 2024 a mis fin à cette organisation par
actions.

\cvrubric{Organisation de journées scientifiques}

Dans le cadre du \textsc{gdr-iasis}, j'ai organisé deux journées thématiques
--- contact avec les intervenants, logistique, communication, animation ---
réunissant chacune une centaine de participants.

\begin{cvfields}
  \item[Janvier 2023] Apprentissage faiblement supervisé en télédétection,
    co-labellisée \textsc{comet cnes}
  \item[Décembre 2024] Télédétection et climat, co-labellisée \textsc{comet
    cnes} et \textsc{gdr} Climat
\end{cvfields}

\cvrubric{Responsabilités au sein de conférences internationales}

\cventry[2027]{Local chair de \textsc{ssp}~2027}{\textsc{ieee} Statistical Signal
  Processing Workshop, Annecy}
Point de contact principal avec le centre de congrès d'Annecy pour l'ensemble
des aspects logistiques --- réservation des salles, gestion des espaces,
coordination technique. J'ai également établi le budget prévisionnel initial
avec le comité d'organisation international.

\cventry[2027]{Équipe de communication d'\textsc{eusipco}~2027}{}

\cventry[2023]{Social media chair d'\textsc{eusipco}~2023}{}


\subsection{Diffusion et valorisation}
\label{subsec:cv-diffusion}

\cvrubric{Activités éditoriales et expertise}

\begin{cvfields}
  \item[Édition] Éditeur associé du \emph{Springer Journal of Applied Signal
    Processing} (\textsc{jasp}) de janvier à décembre 2025. J'ai mis fin à
    cette fonction après le changement de politique éditoriale de Springer
    intervenu en décembre 2025.
  \item[Revues] Relecteur régulier --- une dizaine d'articles par an --- pour
    \emph{\textsc{ieee} Transactions on Signal Processing},
    \emph{\textsc{ieee} Journal of Selected Topics in Signal Processing},
    \emph{\textsc{ieee} Signal Processing Letters}, \emph{Signal Processing}
    (Elsevier), \emph{\textsc{ieee} Transactions on Geoscience and Remote
    Sensing}, \emph{\textsc{ieee} Journal of Selected Topics in Applied Earth
    Observations and Remote Sensing} et \emph{\textsc{ieee} Geoscience and
    Remote Sensing Letters}.
  \item[Conférences] Relecteur pour \textsc{eusipco}, \textsc{icassp},
    \textsc{gretsi} et \textsc{ssp}.
  \item[Expertises] Un projet \textsc{anr}~--- \textsc{astrid}
    (Accompagnement spécifique des travaux de recherches et d'innovation
    défense) et un projet \textsc{anr} Bienvenüe (programme régional Bretagne).
\end{cvfields}

\cvrubric{Sessions spéciales et tutoriaux}

\cventry[2024]{Graph Learning in Signal Processing \& Machine Learning}{Tutoriel de
  3~h à \textsc{eusipco}~2024, avec A.~Breloy, F.~Bouchard et
  A.~Hippert-Ferrer}

\cventry[2023]{Géométrie riemannienne}{Session spéciale \textsc{ss}12 à
  \textsc{eusipco}~2023 --- organisateur}
J'ai par ailleurs animé la session « Traitements multi-capteurs » au
\textsc{gretsi}~2023.

\cventry[2022]{Riemannian and Information Geometry in Signal Processing and Machine
  Learning}{Tutoriel de 6~h à \textsc{eusipco}~2022, avec A.~Breloy et
  F.~Bouchard}
J'ai également animé la session « Remote Sensing » de la même édition.

\cvrubric{Conférences et séminaires invités}

\cventry[2026]{Apprentissage profond sur matrices de covariance : optimisation
  riemannienne}{Journée \textsc{popills} 2026 ---
  \url{https://popils-day.github.io}}

\cventry[2022]{Riemannian Geometry in Machine Learning}{Laboratoire \textsc{cristal},
  université de Lille, février 2022}

\cvrubric{Logiciels libres}

Le développement logiciel est le prolongement naturel du travail
méthodologique : il conditionne la reproductibilité des résultats publiés et
leur reprise par la communauté.

\begin{cvfields}
  \item[qanat] \emph{Auteur.} Gestion, archivage et traçabilité d'expériences
    scientifiques. \url{https://github.com/ammarmian/qanat}
  \item[robuststats] \emph{Auteur.} Estimation robuste de statistiques
    d'ordre~2. \url{https://github.com/ammarmian/robuststats}
  \item[pyRiemann] \emph{Contributeur.} Bibliothèque de référence pour le
    traitement de données sur la variété \textsc{SPD}.
    \url{https://github.com/pyRiemann/pyRiemann}
  \item[spdlearn] \emph{Contributeur.} Apprentissage profond sur données
    \textsc{sdp}. \url{https://github.com/bruAristimunha/spd_learn}
\end{cvfields}

\cvrubric{Vulgarisation}

Fête de la science 2025 : animation d'un atelier grand public sur
l'intelligence artificielle appliquée à la détection d'avalanches, illustrant
l'apport des algorithmes d'apprentissage à la sécurité en montagne.


\subsection{Participation à des jurys}
\label{subsec:cv-jurys}

\begin{cvfields}
  \item[Jurys de thèse] Je n'ai à ce jour siégé dans aucun jury de thèse,
    l'habilitation en étant la condition.
  \item[Comités de suivi] Comité de suivi individuel de Huy Nguyen,
    laboratoire \textsc{sondra}.
\end{cvfields}

\cvrubric{Comités de sélection}

J'ai siégé dans cinq comités de sélection de maîtres de conférences, dans trois
établissements et sur les deux sections dont relèvent mes activités.

\begin{cvfields}
  \item[\textsc{listic}] Université Savoie Mont-Blanc : un comités en 2023,
    en section \textsc{cnu}~61, et un comité en 2026 en section
    \textsc{cnu}~27.
  \item[\textsc{leme}] \textsc{iut} de Ville-d'Avray : deux comités, en 2023
    et en 2024, tous deux en section \textsc{cnu}~61.
  \item[\textsc{cedric}] \textsc{cnam} : un comité en 2026, en section
    \textsc{cnu}~27.
\end{cvfields}


% ============================================================================
\section{Activités d'enseignement}
\label{sec:cv-enseignement}
% ============================================================================

\subsection{Synthèse}
\label{subsec:cv-enseignement-synthese}

J'enseigne à Polytech Annecy-Chambéry sur l'ensemble du cursus d'ingénieur ---
du cycle préparatoire intégré au master~2 --- ainsi que dans la formation
doctorale de l'école \textsc{sie}. Mes enseignements se répartissent entre
l'informatique (programmation, algorithmique, structures de données, bases de
données, développement web), les mathématiques appliquées (optimisation) et
l'intelligence artificielle. Le détail annuel du service figure en
\cref{sec:cv-annexe-enseignements}.


\subsection{Modules et responsabilités}
\label{subsec:cv-modules}

Je suis responsable pédagogique des modules ci-dessous : montage du cours,
réalisation des supports, coordination de l'équipe d'intervenants, sujets
d'examen et bilan.

\cventry[L1]{\textsc{info}101 --- Introduction à la programmation}{Tronc commun,
  80 étudiants --- \textsc{cm}, \textsc{td}, \textsc{tp}}
Bases de la programmation en Python. Responsable sur le site d'Annecy, en
coordination avec le site du Bourget ; équipe de 4 intervenants.

\cventry[L3]{\textsc{info}501 --- Numération et algorithmique}{Tronc commun,
  200 étudiants --- \textsc{cm}, \textsc{td}, \textsc{tp}}
Représentation des nombres et algorithmique en Python. Responsable d'une équipe
de 9 à 10 intervenants.

\cventry[L3]{\textsc{info}503 --- Algorithmes et structures de données}{Filières
  \textsc{sni} et \textsc{idu} --- \textsc{cm}, \textsc{td}, \textsc{tp}}
Algorithmique en langage C et analyse de structures de données.

\cventry[L3]{\textsc{data}601 --- Enjeux de l'intelligence artificielle}{Tronc
  commun --- \textsc{cm}, \textsc{td}, \textsc{tp}}
Mécanismes des \textsc{ia} génératives, enjeux éthiques et environnementaux.
Montage du cours pour la rentrée 2026.

\cventry[L3]{\textsc{math}641 --- Mathématiques spécialisées : optimisation}{Filière
  \textsc{sni} --- \textsc{cm}, \textsc{td}}
Co-responsable avec Abdourahmane Atto.

\cventry[M1]{\textsc{info}734 --- Développement full-stack}{Filière \textsc{idu}
  --- \textsc{cm}, \textsc{td}, projet}
Application JavaScript complète : front-end, back-end, base de données.

\cventry[M1]{\textsc{data}832 --- Machine learning~I}{Filière \textsc{idu} ---
  \textsc{cm}, \textsc{td}, \textsc{tp}}
Introduction à l'apprentissage automatique, orientée pratique.

\cventry[M2]{\textsc{data}931 --- Théorie de l'apprentissage}{Filière \textsc{idu}
  --- \textsc{cm}, \textsc{td}, \textsc{tp}}
Concepts fondateurs de l'apprentissage automatique et de l'optimisation.

\cventry[M2]{\textsc{isoc}931 --- Introduction à la recherche}{Filière \textsc{idu}
  et école doctorale \textsc{sie} --- \textsc{cm}, \textsc{td}, \textsc{tp}}
Méthodologie de la recherche en informatique, état de l'art, rédaction
scientifique.

\cventry[Doctorat]{Numerical Optimisation: Theory \& Applications}{Formation doctorale de
  l'école \textsc{sie} --- \textsc{cm}, \textsc{td}, \textsc{tp}}
Outils mathématiques des algorithmes d'optimisation continue, application aux
problèmes inverses et à l'apprentissage automatique.

\medskip
\noindent J'interviens par ailleurs dans les modules de modélisation
stochastique et d'apprentissage non supervisé placés sous la responsabilité de
collègues, et j'assure chaque année le suivi de stages pour les filières
\textsc{idu} et \textsc{sni}. J'ai enfin été responsable des semestres~7 et~8
de la filière \textsc{sni} de septembre 2023 à août 2026.


% ============================================================================
\section{Responsabilités collectives}
\label{sec:cv-responsabilites}
% ============================================================================

\cventry[depuis 2023]{Chargé de mission « intelligence artificielle dans les
  formations »}{Polytech Annecy-Chambéry --- 12~h équivalent \textsc{td}}
Référent \textsc{ia} de l'établissement, je définis la stratégie pédagogique en
matière d'intelligence artificielle et coordonne les enseignements associés dans
les différentes filières.

\begin{cvfields}
  \item[Cadrage] Animation de groupes de réflexion sur l'usage de
    l'\textsc{ia} par les élèves
  \item[Formation] Accompagnement des collègues souhaitant intégrer
    l'\textsc{ia} générative à leurs cours
  \item[Plateforme] Développement d'une plateforme d'\textsc{ia} souveraine
    hébergée par l'école, avec encadrement de l'ingénieur de recherche dédié
  \item[Communication] Diffusion de ces enjeux auprès de l'école et des
    laboratoires de l'université
\end{cvfields}

\cventry[depuis 2024]{Responsable du thème AFUTE}{\textsc{listic} ---
  17 permanents, 8 doctorants}
Voir \cref{subsec:cv-animation}.

\cventry[2023--2026]{Responsable des semestres~7 et~8}{Filière \textsc{sni}, Polytech
  Annecy-Chambéry}
De septembre 2023 à août 2026.



% ============================================================================
\section{Production scientifique}
\label{sec:cv-production}
% ============================================================================

Ma production compte à ce jour 11~articles de revue internationale à comité de
lecture, 3~chapitres d'ouvrage, 2~communications dans des conférences
d'apprentissage automatique de rang \textsc{a} ou \textsc{a}$^\ast$,
14~communications dans des conférences internationales de traitement du signal
et 7~communications nationales. L'ensemble est déposé sur \textsc{hal}
(extraction disponible ici : \url{https://ammarmian.fr/Publications.html}).

Les noms \underline{soulignés} sont ceux d'étudiants ou de post-doctorants que
j'ai encadrés. La clé portée en regard de chaque référence est celle qui la
désigne dans le reste du mémoire. Chaque catégorie est donnée par date de
parution décroissante.

\cvrubric{Revues internationales à comité de lecture (11)}

\begin{cvpublist}
  \cvpub{SP26}
  \cvpub{JOE26}
  \cvpub{JSTARS25}
  \cvpub{SP24}
  \cvpub{JOE24}
  \cvpub{TGRS23}
  \cvpub{JSTARS20}
  \cvpub{SP20}
  \cvpub{TSP19}
  \cvpub{TGRS19}
  \cvpub{GRSL19}
\end{cvpublist}

\cvrubric{Chapitres d'ouvrage (3)}

\begin{cvpublist}
  \cvpub{CHAP22}
  \cvpub{CHAP21a}
  \cvpub{CHAP21b}
\end{cvpublist}

\cvrubric{Conférences d'apprentissage automatique (2)}

\begin{cvpublist}
  \cvpub{ICML24}
  \cvpub{ECML23}
\end{cvpublist}

\cvrubric{Conférences internationales de traitement du signal (14)}

\begin{cvpublist}
  \cvpub{EUSIPCO26a}
  \cvpub{EUSIPCO26b}
  \cvpub{IGARSS25}
  \cvpub{EUSIPCO25}
  \cvpub{MEAS25}
  \cvpub{ICASSP24}
  \cvpub{IGARSS24}
  \cvpub{EUSIPCO22}
  \cvpub{IGARSS22}
  \cvpub{EUSIPCO21}
  \cvpub{ICASSP19}
  \cvpub{ICASSP18}
  \cvpub{EUSIPCO18}
  \cvpub{EUSIPCO17}
\end{cvpublist}

\cvrubric{Conférences nationales (7)}

\begin{cvpublist}
  \cvpub{RSV25}
  \cvpub{GRETSI23a}
  \cvpub{GRETSI23b}
  \cvpub{GRETSI23c}
  \cvpub{GRETSI22}
  \cvpub{GRETSI19}
  \cvpub{GRETSI17}
\end{cvpublist}

\cvrubric{Thèse de doctorat}

\begin{cvpublist}
  \cvpub{THESE19}
\end{cvpublist}


% ============================================================================
\section{Annexe : service d'enseignement détaillé}
\label{sec:cv-annexe-enseignements}
% ============================================================================

Sauf mention contraire, les enseignements ci-dessous sont dispensés à Polytech
Annecy-Chambéry en formation initiale. Les volumes sont donnés en heures
effectives.

\cvrubric{Année 2020--2021}

\begin{tabularx}{\textwidth}{Xccrrrr}
  \toprule
  \tableheadline{Intitulé} & \tableheadline{Niv.} &
  \tableheadline{Form.} & \tableheadline{Eff.} &
  \tableheadline{\textsc{cm}} & \tableheadline{\textsc{td}} &
  \tableheadline{\textsc{tp}} \\
  \midrule
  Algorithmique\footnote{Enseignement dispensé à la faculté \textsc{scem}
    (L1 \textsc{sts} \textsc{miashs}-éco), hors Polytech Annecy-Chambéry.} & L1 & \textsc{fi} &     &     &     & 8     \\
  \textsc{info}501 --- Numération et algorithmique     & L3 & \textsc{fi} & 200 &     & 21,50 & 21,33 \\
  \textsc{info}502 --- Bases de données                & L3 & \textsc{fi} & 24  &     &       & 24    \\
  \textsc{info}890 --- Algorithmique et bases de données\footnote{Enseignement
    dispensé en formation continue (\textsc{ige}4 \textsc{mp ap}, cohorte
    apprentissage), hors formation initiale.} & M1 & \textsc{fc} &     &     & 12    &       \\
  \midrule
  \textbf{Total} & & & & \textbf{0} & \textbf{33,50} & \textbf{53,33} \\
  \bottomrule
\end{tabularx}

\cvrubric{Année 2021--2022}

\begin{tabularx}{\textwidth}{Xccrrrr}
  \toprule
  \tableheadline{Intitulé} & \tableheadline{Niv.} &
  \tableheadline{Form.} & \tableheadline{Eff.} &
  \tableheadline{\textsc{cm}} & \tableheadline{\textsc{td}} &
  \tableheadline{\textsc{tp}} \\
  \midrule
  \textsc{info}101 --- Introduction à la programmation & L1 & \textsc{fi} & 80  &      &     & 8    \\
  \textsc{info}501 --- Numération et algorithmique     & L3 & \textsc{fi} & 200 & 21   & 25  &      \\
  \textsc{info}502 --- Bases de données                & L3 & \textsc{fi} & 24  &      & 9   & 8    \\
  \textsc{data}931 --- Machine learning                & M2 & \textsc{fi} & 12  & 4,50 & 3   & 2    \\
  \midrule
  \textbf{Total} & & & & \textbf{25,50} & \textbf{37} & \textbf{18} \\
  \bottomrule
\end{tabularx}

\cvrubric{Année 2022--2023}

\begin{tabularx}{\textwidth}{Xccrrrr}
  \toprule
  \tableheadline{Intitulé} & \tableheadline{Niv.} &
  \tableheadline{Form.} & \tableheadline{Eff.} &
  \tableheadline{\textsc{cm}} & \tableheadline{\textsc{td}} &
  \tableheadline{\textsc{tp}} \\
  \midrule
  \textsc{info}101 --- Introduction à la programmation & L1 & \textsc{fi} & 80  & 6   & 9  & 24 \\
  \textsc{info}501 --- Numération et algorithmique     & L3 & \textsc{fi} & 200 & 21  & 25 & 24 \\
  \textsc{info}502 --- Bases de données                & L3 & \textsc{fi} & 24  &     & 9  & 24 \\
  \textsc{info}701 --- Programmation Python (\textsc{iae} Annecy) & L3 & \textsc{fi} & 24 & 15 & & \\
  \textsc{data}731 --- Modélisation stochastique       & M1 & \textsc{fi} & 24  & 4,5 &    &    \\
  \textsc{data}832 --- Théorie de l'apprentissage      & M1 & \textsc{fi} & 24  & 3   & 3  & 4  \\
  \textsc{data}931 --- Machine learning                & M2 & \textsc{fi} & 12  & 3   & 3  & 4  \\
  \textsc{info}942 --- Apprentissage non supervisé     & M2 & \textsc{fi} & 24  & 4,5 &    &    \\
  \textsc{isoc}931 --- Introduction à la recherche     & M2 & \textsc{fi} & 12  & 3   & 3  & 4  \\
  Suivi de stage                                       & M2 & \textsc{fi} & 4   &     & 12 &    \\
  \midrule
  \textbf{Total} & & & & \textbf{60} & \textbf{64} & \textbf{84} \\
  \bottomrule
\end{tabularx}

\cvrubric{Année 2023--2024}

\begin{tabularx}{\textwidth}{Xccrrrr}
  \toprule
  \tableheadline{Intitulé} & \tableheadline{Niv.} &
  \tableheadline{Form.} & \tableheadline{Eff.} &
  \tableheadline{\textsc{cm}} & \tableheadline{\textsc{td}} &
  \tableheadline{\textsc{tp}} \\
  \midrule
  \textsc{info}101 --- Introduction à la programmation & L1 & \textsc{fi} & 80  & 4,5 &     &    \\
  \textsc{info}501 --- Numération et algorithmique     & L3 & \textsc{fi} & 200 & 21  & 9   & 48 \\
  \textsc{info}502 --- Bases de données                & L3 & \textsc{fi} & 24  &     & 9   & 12 \\
  \textsc{math}642 --- Optimisation                    & L3 & \textsc{fi} & 12  & 4,5 & 3   &    \\
  \textsc{data}731 --- Modélisation stochastique       & M1 & \textsc{fi} & 24  & 1,5 &     & 8  \\
  \textsc{info}734 --- Développement full-stack        & M1 & \textsc{fi} & 24  & 12  & 24  &    \\
  \textsc{data}832 --- Théorie de l'apprentissage      & M1 & \textsc{fi} & 24  & 3   & 3   & 4  \\
  \textsc{data}931 --- Machine learning                & M2 & \textsc{fi} & 12  & 3   & 3   & 4  \\
  \textsc{isoc}931 --- Introduction à la recherche     & M2 & \textsc{fi} & 12  & 1,5 & 4,5 & 10 \\
  Suivi de stage                                       & M2 & \textsc{fi} & 3   &     & 9   &    \\
  \midrule
  \textbf{Total} & & & & \textbf{51} & \textbf{64,5} & \textbf{86} \\
  \bottomrule
\end{tabularx}

\cvrubric{Année 2024--2025}

\begin{tabularx}{\textwidth}{Xccrrrr}
  \toprule
  \tableheadline{Intitulé} & \tableheadline{Niv.} &
  \tableheadline{Form.} & \tableheadline{Eff.} &
  \tableheadline{\textsc{cm}} & \tableheadline{\textsc{td}} &
  \tableheadline{\textsc{tp}} \\
  \midrule
  \textsc{info}101 --- Introduction à la programmation & L1   & \textsc{fi} & 80  & 6    & 10,5 & 12 \\
  \textsc{info}201 --- Systèmes d'exploitation         & L1   & \textsc{fi} & 15  &      &      & 15 \\
  \textsc{info}501 --- Numération et algorithmique     & L3   & \textsc{fi} & 200 & 36   & 21   & 36 \\
  \textsc{info}502 --- Bases de données                & L3   & \textsc{fi} & 24  &      & 4,5  & 12 \\
  \textsc{math}642 --- Optimisation                    & L3   & \textsc{fi} & 12  & 4,5  & 3    &    \\
  \textsc{data}731 --- Modélisation stochastique       & M1   & \textsc{fi} & 24  & 1,5  &      & 8  \\
  \textsc{info}734 --- Développement full-stack        & M1   & \textsc{fi} & 24  & 12   & 24   &    \\
  \textsc{data}832 --- Théorie de l'apprentissage      & M1   & \textsc{fi} & 24  & 9    & 6    & 8  \\
  \textsc{data}931 --- Machine learning                & M2   & \textsc{fi} & 12  & 3    & 3    & 4  \\
  \textsc{isoc}931 --- Introduction à la recherche     & M2   & \textsc{fi} & 12  & 1,5  & 7,5  & 10 \\
  \textsc{pr\&d sni}                                   & M2   & \textsc{fi} & 2   &      & 15   &    \\
  Numerical Optimisation: Theory \& Applications       & Doc. & \textsc{fi} & 20  & 13,5 & 6    & 4  \\
  Suivi de stage                                       & M2   & \textsc{fi} &     &      & 24   &    \\
  \midrule
  \textbf{Total} & & & & \textbf{87} & \textbf{124,5} & \textbf{109} \\
  \bottomrule
\end{tabularx}


% ============================================================================
\section{Annexe : fiches d'encadrement doctoral}
\label{sec:cv-annexe-encadrements}
% ============================================================================

\cventry[2021--2024]{Olivier Lerda}{Détection robuste pour un sonar à croix de Mills}
\begin{cvfields}
  \item[Période] Début en avril 2021, soutenance le 17 décembre 2024
  \item[Financement] Ingénieur chez iXBlue, thèse préparée à mi-temps en
    parallèle de son poste --- il ne s'agit pas d'un contrat \textsc{cifre}
  \item[Encadrement] 33~\%, avec Guillaume Ginolhac et Jean-Philippe Ovarlez
  \item[Discipline] \textsc{stic} --- traitement de l'information
  \item[École doctorale] Sciences, ingénierie, environnement (Chambéry)
  \item[Laboratoire] \textsc{listic}, Annecy --- soutenance à Chambéry
  \item[Jury] Nicolas Le Bihan (président) ; Philippe Forster et Éric
    Chaumette (rapporteurs) ; Angélique Drémeau (examinatrice)
  \item[\textsc{doi}] \href{https://doi.org/10.70675/342ac9b9z6383z42acza2aez69f69e3bbfcf}%
    {10.70675/342ac9b9z6383z42acza2aez69f69e3bbfcf}
  \item[Publications] \cvcite{JOE24,JOE26,GRETSI23a}
  \item[Devenir] Ingénieur de recherche chez Exail
\end{cvfields}

\cventry[2022--2025]{Douba Jafuno}{Classification de signaux \textsc{gpr}}
\begin{cvfields}
  \item[Période] Début le 1\textsuperscript{er} juin 2022, soutenance le
    4 novembre 2025
  \item[Financement] Contrat \textsc{cifre} avec Géolithe
  \item[Encadrement] 45~\%, avec Guillaume Ginolhac et Nickolas Stelzenmuller
  \item[Discipline] \textsc{stic} --- traitement de l'information
  \item[École doctorale] Sciences, ingénierie, environnement (Chambéry)
  \item[Laboratoire] \textsc{listic}, Annecy --- soutenance à Chambéry
  \item[Jury] Thomas Oberlin (président) ; Sylvain Chevallier et Audrey
    Giremus (rapporteurs) ; Sophie Giffard-Roisin (examinatrice)
  \item[\textsc{doi}] \href{https://doi.org/10.70675/2b4df77cz09d2z4ad1z8c43z5ab013c30172}%
    {10.70675/2b4df77cz09d2z4ad1z8c43z5ab013c30172}
  \item[Publications] \cvcite{JSTARS25,IGARSS24,GRETSI23b}
\end{cvfields}

\cventry[depuis 2024]{Matthieu Verlynde}{Apprentissage frugal pour la télédétection}
\begin{cvfields}
  \item[Période] Début le 1\textsuperscript{er} octobre 2024, soutenance
    prévue en 2027
  \item[Financement] Bourse ministérielle
  \item[Encadrement] 60~\%, avec Yajing Yan
  \item[École doctorale] \textsc{sie}
  \item[Publications] \cvcite{IGARSS25,EUSIPCO25}
\end{cvfields}

\cventry[depuis 2025]{Wenli Shi}{Réseaux de diffusion parcimonieux sur données non
  textuelles}
\begin{cvfields}
  \item[Période] Début en décembre 2025, soutenance prévue en 2028
  \item[Financement] Contrat \textsc{cifre} avec Givaudan
  \item[Encadrement] Co-direction à 33~\%, avec Frédéric Pascal (\textsc{l2s},
    université Paris-Saclay) et Matthieu Jonckheere (\textsc{dr cnrs},
    \textsc{laas})
  \item[École doctorale] STIC, Paris-Saclay
\end{cvfields}

\cventry[à partir de 2026]{Roman Wild}{Phase géométrique en \textsc{polinsar}}
\begin{cvfields}
  \item[Période] Début Octobre 2026
  \item[Financement] Chaire \textsc{miai}
  \item[Encadrement] 33~\%, avec Nicolas Le Bihan et Pierre-Olivier Amblard
  \item[Laboratoire] \textsc{gipsa}-lab, équipe \textsc{gaia}
\end{cvfields}
